\clearpage
\renewcommand{\sectioncolor}{slate}
\section{Appendix: complete inventory of Part VI}
\markboth{Appendix}{}

\subsection{Every metaphor and blend, with its page}

\begin{center}\small
\begin{tabular}{@{}>{\raggedright\arraybackslash}p{6.4cm}>{\raggedright\arraybackslash}p{1.6cm}>{\raggedright\arraybackslash}p{1.6cm}>{\raggedright\arraybackslash}p{6.5cm}@{}}
\toprule
\rowcolor{slate!14}
\textbf{Name} & \textbf{book p.} & \textbf{pdf p.} & \textbf{What it buys}\\
\midrule
\multicolumn{4}{@{}l}{\textbf{\color{csA}Case Study 1}}\\
\blend{Cartesian Plane blend}            & 385--386 & 404--405 & geometry $\leftrightarrow$ algebra; $n$-dimensional visualisation\\
\meta{Functions Are Numbers}             & 386      & 405      & arithmetic on functions; later, power series\\
\blend{Unit Circle blend} (4 domains)    & 388--392 & 407--411 & arcs acquire length; $a,b$ acquire numbers\\
\meta{Trigonometry metaphor}             & 393      & 412      & angles become numbers; $\cos\pi=-1$\\
\meta{Recurrence Is Circularity}         & 395      & 414      & periodicity; the second meaning of $\pi$\\
\blend{Trigonometry blend}               & 396      & 415      & sine and cosine as \emph{functions}\\
\meta{Polar--Trigonometry metaphor}      & 396      & 415      & $(r,\theta)$ coordinates\\
\midrule
\multicolumn{4}{@{}l}{\textbf{\color{csB}Case Study 2}}\\
\meta{Multiplication Is Addition}        & 405      & 424      & the exponential, as arithmetized metaphor\\
\blend{Exponential blend}                & 406      & 425      & the log/exp one-to-one correspondence\\
\meta{Time Is Distance}                  & 407      & 426      & step 1 of mathematicizing change\\
\meta{Logical Independence Is Geometrical Orthogonality} & 407 & 426 & the two axes\\
\meta{Instantaneous Speed Is Average Speed over an Infinitely Small Interval} & 408 & 427 & the derivative\\
\meta{Change Is Motion}                  & 408      & 427      & rate of change in general\\
\meta{Numbers Are Wholes That Are Sums of Their Parts} & 412 & 431 & numbers as sums\\
\meta{Infinite Sums Are Limits of Sequences of Partial Sums} & 412 & 431 & BMI special case\\
\blend{Power-Series blend} (6 domains)   & 413      & 432      & term-by-term differentiation\\
\midrule
\multicolumn{4}{@{}l}{\textbf{\color{csC}Case Study 3}}\\
\meta{Fundamental Metonymy of Algebra}   & 421      & 440      & solving for an unknown ``number''\\
\meta{Numbers Are Points on a Line}      & 424      & 443      & the Number-Line blend\\
\meta{Multiplication by $-1$ Is Rotation to the Symmetry Point} & 424--425 & 443--444 & negatives as rotation\\
\blend{Multiplication-Rotation blend}    & 425      & 444      & arithmetizes $R_{-1}$\\
\blend{Rotation-Plane blend}             & 426      & 445      & vector addition, scalar multiplication\\
\meta{$90^{\circ}$ Rotation metaphor}    & 428--429 & 447--448 & $n^{2}=-1$, hence $i$; the complex plane\\
\midrule
\multicolumn{4}{@{}l}{\textbf{\color{csD}Case Study 4}}\\
\blend{Trigonometric Complex Plane blend}& 437      & 456      & $\cis$; complex multiplication in polar form\\
\meta{A Function Is a Set of Ordered Pairs} & 444   & 463      & what ``$=$'' asserts between two rules\\
\blend{Taylor Series blend}              & 443--444 & 462--463 & $a_n=f^{(n)}(0)/n!$\\
\blend{Euler's blend} (5 domains)        & 447      & 466      & the whole identity\\
\bottomrule
\end{tabular}
\end{center}

\subsection{Navigation}

\begin{center}\small
\begin{tabular}{@{}>{\raggedright\arraybackslash}p{7.6cm} c c c@{}}
\toprule
\rowcolor{slate!14}
\textbf{Section} & \textbf{book pp.} & \textbf{pdf pp.} & \textbf{offset}\\
\midrule
Part VI opening --- ``$e^{\pi i}+1=0$''                        & 381--382 & 400--401 & \multirow{5}{*}{$+19$}\\
Case Study 1 · Analytic Geometry and Trigonometry               & 383--398 & 402--417 & \\
Case Study 2 · What Is $e$?                                     & 399--419 & 418--438 & \\
Case Study 3 · What Is $i$?                                     & 420--432 & 439--451 & \\
Case Study 4 · $e^{\pi i}+1=0$                                  & 433--452 & 452--471 & \\
\midrule
References                                                      & 453--472 & 472--491 & \\
Index                                                           & 473--     & 492--    & \\
\bottomrule
\end{tabular}
\end{center}
\vspace{-2pt}
\begin{center}\begin{minipage}{0.93\textwidth}\footnotesize\color{slate}
\textbf{Table 8.}\; Page correspondence for
\texttt{Where\_Mathematics\_Come\_From\_\_by\_G.Lakoff\_R.Nunez\_2001.pdf} (514 pages).
Add \textbf{19} to a printed book page to get the PDF page.
\end{minipage}\end{center}

\vspace{22pt}
\begin{center}
\begin{tikzpicture}
\node[rounded corners=5pt,fill=deepblue,text=white,inner sep=13pt,text width=15.6cm,align=center]
 {\mbox{\large\bfseries $e^{\pi i}+1=0$}\\[7pt]
  \normalfont\small It brings together $e$ (change), $\pi$ (recurrence), $i$ (rotation),
  $1$ and $0$ (the two identities) --- and it is true because, under the metaphors and blends of
  classical mathematics, \textbf{that is what those symbols mean}.};
\end{tikzpicture}
\end{center}
\vspace{10pt}
